\begin{claim}{F.30}{Proved}{Finite affine product-one images when the first scalar generates the scalar-field extension}
Use the actual bases $Y_i^-=X_i$ and
$Y_i^+=\mathcal B_+(X_i)$ under the based
isomorphism of F.23.13.
For a finite field $k$ and continuous
$\rho:V_\pm\to\operatorname{Aff}_1(k)$, let
$\chi$ be its scalar character and put
$C=\chi(R_\pm)$,
$k_1=\mathbb F_p(C)$,
$k_0=\mathbb F_p(\chi(V_\pm))$ and
$r=\chi(Y_1^\pm)$.
The exact locus
\[
             \prod_{i=1}^7\chi(Y_i^\pm)=1,
                    \qquad k_0=k_1(r)
                                                               \tag{F.30.1}
\]
has the image equality
\[
                         \rho(\mathcal K_\pm)=\rho(R_\pm).
                                                               \tag{F.30.2}
\]
All characteristics, finite extension degrees
and larger original coefficient fields are included.
On the part $C\ne\{1\}$, put
$q_1=|k_1|$, $m=[k_0:k_1]$ and
$\chi_j=\chi^{q_1^j}$, $0\leq j<m$.
The summed continuous $H^1$ restriction for
these $m$ characters is injective. Its cocycle
kernel consists exactly of the principal tuples
whose parameters sum to zero, and its rank is
$6m+1$ on $7m$ coordinates.
\end{claim}

\begin{proof}
\emph{Actual fields and the completed scalar cases.}
Use the negative groups $K=\mathcal K_-\leq R=R_-\leq V=V_-$
first. Write their product-one character as
$(r,s,t,u,v,w,c^{-1})$, where $c=rstuvw$.
The normal-generation equality of F.21 maps
through the actual marginal surjection of F.22
to $\overline{\langle\!\langle K\rangle\!\rangle}_R=R$.
Its image under the continuous abelian scalar
map gives $\chi(K)=\chi(R)$.
The actual deleted generators $X_4,\ldots,X_7$
and (F.30.1) therefore give
\[
 \begin{gathered}
 C=\chi(K)=\chi(R)=\langle u,v,w,c^{-1}\rangle,
 \qquad k_1=\mathbb F_p(u,v,w,c),\\
 k_0=\mathbb F_p(r,s,t,u,v,w)=k_1(r).
 \end{gathered}                                                \tag{F.30.3}
\]
All six parameters are nonzero.
For $C=\{1\}$, the image of $R$ is an abelian
translation group, so the same actual
normal-generation map proves (F.30.2).
For $C\ne\{1\}$ and $k_0=k_1$, F.25 gives
the image equality. F.24 supplies the individual
restriction injection and the stated cocycle
kernel with $m=1$.

We now treat $C\ne\{1\}$ and $k_1\subsetneq k_0$.
Equation (F.30.3) implies $r\notin k_1$, hence
$r-1\ne0$. This gives an actual retained-generator
principal pivot, independently of the value of $c$.

\emph{Exact coefficient and principal-quotient maps.}
The word derivations use the coefficient homomorphism
\[
 \begin{aligned}
 \mathbb Z[t_1^{\pm1},\ldots,t_7^{\pm1}]
    &\longrightarrow A=
       \mathbb Z[r^{\pm1},s^{\pm1},t^{\pm1},u^{\pm1},v^{\pm1},w^{\pm1}],\\
 (t_1,t_2,t_3,t_4,t_5,t_6,t_7)
    &\longmapsto(r,s,t,u,v,w,(rstuvw)^{-1}).
 \end{aligned}                                                \tag{F.30.4}
\]
Thus an exponent vector $(a_i)$ maps to
$r^{a_1-a_7}s^{a_2-a_7}t^{a_3-a_7}
 u^{a_4-a_7}v^{a_5-a_7}w^{a_6-a_7}$.
These are derivations on the specified discrete
free words. Their evaluations in a finite field
give the actual cocycle rule.
Every seven-tuple of generator values defines
a homomorphism to the finite affine group
and extends uniquely to $V=\widehat F_7$.
Consequently the continuous cocycle space is
exactly $k_0^7$.

Let $\delta=(r,s,t,u,v,w,c^{-1})^{\mathsf T}-\mathbf1$.
Its first coordinate is $\delta_1=r-1\ne0$.
The actual principal quotient uses the maps
\[
 \begin{aligned}
 P_1(z)&=(z_i-\delta_i(r-1)^{-1}z_1)_{i=2}^7,
       &S_1(x)&=(0,x_1,\ldots,x_6)^{\mathsf T},\\
 P_1S_1&=I_6,
       &S_1P_1(z)&=z-\delta(r-1)^{-1}z_1.
 \end{aligned}                                                \tag{F.30.5}
\]
They identify
$Q_\chi=H^1_{\rm cont}(V,k_{0,\chi})$
with six coordinates corresponding, in order,
to the original $X_2,\ldots,X_7$.
The universal maps are defined over $A[(r-1)^{-1}]$;
the proved nonvanishing justifies every
finite-field use of that localization.

Let $P_0=X_1\cdots X_7$ be the actual boundary word.
Its character is one, and its derivative and
induced quotient covector are
\[
 \begin{gathered}
 D(P_0)=(1,r,rs,rst,rstu,rstuv,c),\\
 b=D(P_0)S_1=(r,rs,rst,rstu,rstuv,c),\\
 D(P_0)\delta=0,\qquad D(P_0)=bP_1.
 \end{gathered}                                                \tag{F.30.6}
\]
The vanishing is the telescoping identity
$D(P_0)\delta=\chi(P_0)-1$; substituting
(F.30.5) proves the last displayed equality.

For $\beta=(s,st,stu,stuv,stuvw)$, define
\[
 \begin{aligned}
 \pi:Q_\chi&\longrightarrow k_0,
       &\pi(z)&=bz,\\
 \iota:k_0^5&\longrightarrow Q_\chi,
       &\iota(x)&=(-\textstyle\sum_{i=1}^5\beta_ix_i,x_1,\ldots,x_5),\\
 \sigma:k_0&\longrightarrow Q_\chi,
       &\sigma(a)&=r^{-1}a e_1.
 \end{aligned}                                                \tag{F.30.7}
\]
One has $\pi\sigma=I_1$,
$\ker\pi=\operatorname{im}\iota=:W_\chi$ and
$\iota(z_2,\ldots,z_6)+\sigma\pi(z)=z$.
In particular the coefficient basis and its
inverse are explicitly
\[
 \begin{gathered}
 J=(e_2-se_1,e_3-ste_1,e_4-stue_1,
                 e_5-stuve_1,e_6-stuvwe_1,r^{-1}e_1),\\
                  J^{-1}(z)=(z_2,z_3,z_4,z_5,z_6,bz).
 \end{gathered}                                                \tag{F.30.8}
\]
The displayed decomposition proves both
basis composites without an assumed invariant complement.

\emph{The actual source operators needed for separation.}
Retain the eleven F.23 automorphisms
$B_l(a)=u_l\alpha_l(a)u_l^{-1}$, $0\leq l\leq10$.
They and their inverses preserve $K$.
Their pure Artin action fixes every abelian
character, and their cocycle matrices $A_l$
therefore satisfy $A_l\delta=\delta$.
Set $\overline A_l=P_1A_lS_1$.

Every pure $\alpha_l$ fixes $P_0$. The cocycle
formula for conjugation gives
$d(B_l(P_0))=\chi(u_l)d(P_0)$ because
$\chi(P_0)=1$. Combining this with (F.30.6)
gives $\pi\overline A_l=\chi(u_l)\pi$.
Thus $W_\chi$ is invariant under all actual
source operators, and $\pi$ is their exact
one-dimensional boundary quotient.

Put $\widetilde\iota=S_1\iota$.
The literal original-word calculation gives
the integral identity
\[
 (r-1)(A_0\widetilde\iota-r\widetilde\iota)
       =\delta\,\operatorname{row}_1(A_0\widetilde\iota).
                                                               \tag{F.30.9}
\]
Apply $P_1$ and cancel the proved nonzero
$r-1$. This gives the actual equality
$\overline A_0\iota=r\iota$ on the five-dimensional
invariant subspace. The same word calculation
gives the boundary entries
\[
 \begin{array}{c|c|c}
 l&\overline A_l|W_\chi&\chi(u_l)\text{ on }Q_\chi/W_\chi\\\hline
 0&rI_5&1\\
 2&*&rt\\
 7&*&t^{-1}
 \end{array}.
                                                               \tag{F.30.10}
\]
The stars stand for the actual core operators;
only the displayed entries are used.
The indices are the original zero-based
$0=\texttt{old0}$,
$2=\texttt{old\_1\_1}$ and
$7=\texttt{old\_-1\_1}$, respectively the
first, third and eighth source positions.
The certificate compares all entries of the
three full seven-coordinate source actions
with direct evaluations of those labelled
decorated Artin words.

\emph{Frobenius transport and all three type comparisons.}
For $0\leq j<m$, put
$r_j=r^{q_1^j}$, $s_j=s^{q_1^j}$ and
$t_j=t^{q_1^j}$. The parameters $u,v,w,c$
are fixed by these powers.
The coefficient map
\[
 d\longmapsto d^{q_1^j},\qquad
 k_{0,\chi}\longrightarrow k_{0,\chi_j}
                                                               \tag{F.30.11}
\]
is additive and semilinear. It preserves the
cocycle rule, principal cocycles and continuity,
and commutes with each actual source precomposition.
The inverse is the power $q_1^{m-j}$, with
the identity at $j=0$.
It transports every displayed inclusion,
quotient, section and inverse basis to
$Q_j=H^1_{\rm cont}(V,k_{0,\chi_j})$ and $W_j$.
All characters agree on $K$, since
$\chi(K)=C\subset k_1^\times$.

The elements $r_0,\ldots,r_{m-1}$ are distinct
and each differs from $1$.
An equality of two would make a nonidentity
$k_1$-Frobenius power fix the generator $r$
and hence all of $k_0=k_1(r)$.
For $0<a<m$, the polynomial $X^{q_1^a}-X$
has degree smaller than $|k_0|$, so cannot
vanish on the whole field.
Also a conjugate equal to $1$ would give
$r=1\in k_1$, contrary to the proper extension.

Take a nonzero source-invariant subspace
$U\leq Q_j$ with no proper nonzero invariant
subspace. A nonzero invariant intersection
with $W_j$ gives $U\leq W_j$.
For zero intersection, $\pi_j|U$ is injective
and identifies $U$ with the one-dimensional
boundary character. These are the two
possibilities through the exact maps (F.30.7).

For distinct $j,l$, the following comparisons
exclude every source-intertwining isomorphism
between such subspaces. Two subspaces in
$W_j,W_l$ have source-$0$ scalars $r_j,r_l$;
an intertwiner $F$ would satisfy
$(r_j-r_l)F=0$, hence $F=0$.
A core subspace and a boundary subspace
have source-$0$ scalars $r_j$ and $1$,
and the same argument excludes an isomorphism.
For two boundary subspaces, source $7$
would force $t_j^{-1}=t_l^{-1}$, so $t_j=t_l$.
Source $2$ would then force
$r_jt_j=r_lt_l$, hence $r_j=r_l$,
contrary to the proved distinctness.
All character values cancelled here are
nonzero. Inverse intertwiners cover both
orientations of the mixed comparison.
These are all three unordered type comparisons.

\emph{Summed restriction and the finite affine image.}
Consider the well-defined actual restriction map
\[
 \Phi:\bigoplus_{j=0}^{m-1}Q_j
      \longrightarrow H^1_{\rm cont}(K,k_{0,\chi|K}),
      \qquad ([d_j])_j\longmapsto\sum_j[d_j|K].
                                                               \tag{F.30.12}
\]
It commutes with every actual source operator.
Every individual restriction is injective
by F.24, because $\chi_j(R)=C\ne\{1\}$.
Under the supposition that $\ker\Phi$ is
nonzero, choose a nonzero invariant subspace
$T$ inside it with smallest dimension.
Each nonzero coordinate projection $p_j|T$
is injective, and its image has no proper
nonzero invariant subspace.
Two nonzero projections give the explicit
mutually inverse intertwiners
\[
 p_l\circ(p_j|T)^{-1}:p_j(T)\longrightarrow p_l(T),
 \qquad p_j\circ(p_l|T)^{-1}:p_l(T)\longrightarrow p_j(T).
                                                               \tag{F.30.13}
\]
The three comparisons exclude this for
$j\ne l$. Hence there is only one nonzero
projection, so $T$ lies in one $Q_j$.
There $\Phi$ is the injective individual
restriction, a contradiction.
This proves the summed $H^1$ injection.

Zero summed restriction of cocycles therefore
forces each cocycle to be principal.
Since all characters agree on $K$ and
$C\ne\{1\}$, the exact kernel and its dimension are
\[
 \left\{\bigl(b_j(\chi_j-1)\bigr)_j:
                      \sum_jb_j=0\right\},\qquad
                     \dim\ker=m-1.
                                                               \tag{F.30.14}
\]
The converse inclusion is direct restriction.
Each principal vector is nonzero, and each
continuous cocycle space has seven generator
coordinates. Thus the stated rank is
$7m-(m-1)=6m+1$.

Apply this result to the tuple
$(d^{q_1^j})_j$ of one cocycle in the trace
kernel of F.26. The cocycle $d$ is principal,
say $d=b(\chi-1)$.
Its Frobenius powers have parameters
$b^{q_1^j}$, uniquely because their principal
vectors are nonzero. Their parameter sum
is exactly the field trace. Consequently
\[
 \operatorname{Tr}_{k_0/k_1}(d|K)=0
 \quad\Longleftrightarrow\quad
 d=b(\chi-1),\qquad \operatorname{Tr}_{k_0/k_1}(b)=0.
                                                               \tag{F.30.15}
\]
The reverse implication follows directly from
$\chi(K)\subset k_1$.
The explicit conjugation and coefficient
homomorphism of F.26 sends every proper
finite affine marginal image to a nonprincipal
element of this trace kernel over $k_0$.
Those maps preserve the scalar character
and both scalar fields. They therefore
preserve the product-one condition and
$k_0=k_1(r)$ of (F.30.1).
Equation (F.30.15) excludes their resulting
nonprincipal element. This proves (F.30.2)
for the negative marginal over every original
finite coefficient field, including fields
larger than $k_0$.

\emph{Integral evidence and actual positive-sign transport.}
The coefficient producer applies precisely
(F.30.4) to the frozen F.24 universal matrices
for the three original sources.
The separate verifier imports neither that
producer nor an algebra package. It uses the
reviewed integer Laurent arithmetic to
reconstruct the three actual decorated Artin
words and compares every entry of their full
seven-coordinate actions with the producer.
Its $45$ exact checks include both inverse
maps in (F.30.8), the boundary kernel and
section, the principal-vector identities,
(F.30.9) and every boundary value in (F.30.10).
The identities have integer coefficients;
the only nonmonomial principal localization
is $r-1$, whose nonvanishing was proved from
the exact field equality before its use.
The evidence manifest binds the statement,
programs, original words, coefficient inputs,
outputs and completed dependencies.

For a positive representation, precomposition
with the actual continuous isomorphism
$\mathcal B_+:V_-\to V_+$ sends the values
at $Y_i^+$ to precisely the negative
coordinates used above.
F.23.13 supplies its inverse and carries
$\mathcal K_-,R_-$ onto their positive
counterparts. Thus it preserves the scalar
images defining $k_0,k_1$, the first scalar
$r$ and all conditions (F.30.1).
Equality of the two negative images gives
equality of the original two positive images
through these specified maps. This proves
(F.30.2) for both signs.
\end{proof}
